% ==================== 7.2 主要创新点 ====================
\section{主要创新点}
\label{sec:innovations}

本文的主要创新点如下：

\textbf{创新点一：提出了面向月面巡视器自主行走的增强五维数字孪生模型}

针对经典数字孪生五维模型在月面巡视器自主行走场景下的局限性，提出了增强的五维模型。该模型创新性地将月面环境纳入物理实体维度，在虚拟实体维度引入自主演化机制，在服务维度增加自主决策支持，在孪生数据维度实现多源异构融合，在连接维度设计时延补偿机制。该模型为月面巡视器自主行走数字孪生系统的构建提供了理论指导，相关成果可推广至其他深空探测器数字孪生系统。

\textbf{创新点二：建立了考虑滑移特性和低重力效应的巡视器动力学数字孪生模型}

针对月面环境下巡视器动力学建模的难题，建立了考虑滑移沉陷、低重力效应的轮-壤相互作用模型和巡视器多体动力学模型。该模型能够准确预测巡视器在月面复杂地形下的运动行为，满足实时性要求。相比传统模型，该模型考虑了月面环境的特殊性，预测精度更高，为巡视器自主行走提供了可靠的动力学支撑。

\textbf{创新点三：提出了基于A-D3QN的混合路径规划算法}

针对月面复杂环境下的路径规划问题，提出了基于注意力机制增强的双深度Q网络（A-D3QN）混合路径规划算法。该算法创新性地将注意力机制引入路径规划网络，使网络能够自适应地关注关键环境因素；设计了全局-局部协同的混合规划策略，实现了规划效率与规划质量的统一。该算法在动态环境下表现出色，为巡视器自主行走提供了有效的规划方法。

以上创新点相互关联、层层递进，形成了面向月面巡视器自主行走的数字孪生技术体系，为解决月面巡视器自主行走难题提供了系统方案。
